% ----------------------------------------------------------
% EPÍGRAFE

%Epígrafe: Elemento opcional e sem título em que o (a) autor (a) apresenta uma citação relacionada ao assunto tratado no trabalho. Deve ser elaborada conforme a ABNT-NBR 10520 (Citações). As citações de até três linhas devem estar entre aspas duplas e as citações com mais de três linhas devem ser destacadas com recuo de 4 cm da margem esquerda, com letra menor que a do texto e sem as aspas.
% % ---------------------------------------------------------
\vspace*{10cm}
\begin{citacao}
Exemplo de \textbf{EPÍGRAFE (opcional)} com mais de 3 linhas: Elemento opcional e sem título em que o (a) autor (a) apresenta uma citação relacionada ao assunto tratado no trabalho. Deve ser elaborada conforme a ABNT-NBR 10520 (Citações). As citações de até três linhas devem estar entre aspas duplas e as citações com mais de três linhas devem ser destacadas com recuo de 4 cm da margem esquerda, com letra menor que a do texto e sem as aspas. (ABNT, 2002).
\end{citacao}

    \vspace*{5cm}
	
		\enquote{Aqui temos um exemplo de como, se você realmente quiser aumentar o tamanho do trabalho e estiver disposto a procurar algo legal, você pode incluir uma \textbf{CITAÇÃO (opcional)} no seu texto. Geralmente você deixa isso de lado para focar em coisas mais urgentes, ou não tem tempo.} (Sobrenome; Nome, 2025).
	
\newpage